\documentclass[11pt]{article}
\usepackage[margin=1in]{geometry}
\usepackage{amsmath, amssymb}
\usepackage{enumitem}
\usepackage{array}
\usepackage{booktabs}
\usepackage{graphicx}
\usepackage{parskip}

\begin{document}

\noindent
\textbf{CS 162 -- Homework 3} \hfill Kaushik Vada \\
Winter 2026 \hfill February 19, 2026

\bigskip
\hrule
\bigskip

%=============================================================================
% PROBLEM 1
%=============================================================================
\section*{1. TLB / Page Table / Cache Possibility Table (10 pts)}

The lookup order is TLB $\to$ Page Table $\to$ Cache.

\renewcommand{\arraystretch}{1.35}
\begin{center}
\begin{tabular}{|c|c|c|c|p{7cm}|}
\hline
\textbf{TLB} & \textbf{Page Table} & \textbf{Cache} & \textbf{Possible?} & \textbf{Under what circumstances?} \\
\hline
Hit & Hit & Hit & Yes & Normal case -- TLB has the translation, page is in memory, and the data is in cache. \\
\hline
Hit & Hit & Miss & Yes & TLB has the translation and the page is valid, but the data isn't in the cache (compulsory/capacity/conflict miss). \\
\hline
Miss & Hit & Hit & Yes & TLB doesn't have the entry (maybe evicted earlier), but the page table has a valid mapping. After getting the physical address, the data is found in cache. \\
\hline
Miss & Hit & Miss & Yes & TLB misses, page table has the mapping, but the data block isn't in cache. Need to go to main memory. \\
\hline
Miss & Miss & Miss & Yes & TLB misses, page table also misses (page fault -- page isn't in physical memory). After handling the page fault and bringing the page from disk, the data won't be in cache either. \\
\hline
Hit & Miss & Hit/Miss & No & A TLB hit means we have a valid translation, which means the page \textit{must} be in physical memory. So the page table can't miss -- TLB hit implies page table hit. \\
\hline
Miss & Miss & Hit & No & If the page table misses, that's a page fault -- the page isn't in physical memory. So there's no way the data is in the cache since cache holds physical memory data. \\
\hline
\end{tabular}
\end{center}

\newpage
%=============================================================================
% PROBLEM 2
%=============================================================================
\section*{2. Data Cache Simulation (50 pts)}

Cache starts with word addresses 0--15. Reference string:
\begin{center}
\texttt{1, 20, 2, 3, 4, 18, 5, 19, 33, 34, 1, 4}
\end{center}

%-----------------------------------------------------------------------------
\subsection*{(a) Direct-Mapped, 1-Word Blocks, 16 Blocks}

Index $= \text{Word Addr} \bmod 16$. Initially block $i$ holds word address $i$.

\medskip
\begin{center}
\begin{tabular}{|c|c|c|l|}
\hline
\textbf{\#} & \textbf{Addr} & \textbf{Index} & \textbf{H/M} \\
\hline
1 & 1 & 1 & Hit \\
2 & 20 & 4 & Miss (had 4, replace w/ 20) \\
3 & 2 & 2 & Hit \\
4 & 3 & 3 & Hit \\
5 & 4 & 4 & Miss (had 20, replace w/ 4) \\
6 & 18 & 2 & Miss (had 2, replace w/ 18) \\
7 & 5 & 5 & Hit \\
8 & 19 & 3 & Miss (had 3, replace w/ 19) \\
9 & 33 & 1 & Miss (had 1, replace w/ 33) \\
10 & 34 & 2 & Miss (had 18, replace w/ 34) \\
11 & 1 & 1 & Miss (had 33, replace w/ 1) \\
12 & 4 & 4 & Hit \\
\hline
\end{tabular}
\end{center}

\medskip
Hits = 5, Misses = 7

\textbf{Hit rate = 5/12 $\approx$ 41.67\%}

\medskip
Final cache state:

\begin{center}
\begin{tabular}{|c|c||c|c|}
\hline
\textbf{Idx} & \textbf{Addr} & \textbf{Idx} & \textbf{Addr} \\
\hline
0 & 0 & 8 & 8 \\
1 & 1 & 9 & 9 \\
2 & 34 & 10 & 10 \\
3 & 19 & 11 & 11 \\
4 & 4 & 12 & 12 \\
5 & 5 & 13 & 13 \\
6 & 6 & 14 & 14 \\
7 & 7 & 15 & 15 \\
\hline
\end{tabular}
\end{center}

\newpage
%-----------------------------------------------------------------------------
\subsection*{(b) Direct-Mapped, 2-Word Blocks, 8 Blocks}

Block \# $= \lfloor \text{Addr} / 2 \rfloor$, \quad Index $= \text{Block \#} \bmod 8$. Each block holds words $\{2k, 2k{+}1\}$.

Initially loaded with addrs 0--15, so indices 0--7 hold blocks 0--7.

\medskip
\begin{center}
\begin{tabular}{|c|c|c|c|l|}
\hline
\textbf{\#} & \textbf{Addr} & \textbf{Blk} & \textbf{Idx} & \textbf{H/M} \\
\hline
1 & 1 & 0 & 0 & Hit (blk 0: \{0,1\}) \\
2 & 20 & 10 & 2 & Miss (had blk 2, replace w/ \{20,21\}) \\
3 & 2 & 1 & 1 & Hit (blk 1: \{2,3\}) \\
4 & 3 & 1 & 1 & Hit (still blk 1) \\
5 & 4 & 2 & 2 & Miss (had blk 10, replace w/ \{4,5\}) \\
6 & 18 & 9 & 1 & Miss (had blk 1, replace w/ \{18,19\}) \\
7 & 5 & 2 & 2 & Hit (blk 2: \{4,5\}) \\
8 & 19 & 9 & 1 & Hit (blk 9: \{18,19\}) \\
9 & 33 & 16 & 0 & Miss (had blk 0, replace w/ \{32,33\}) \\
10 & 34 & 17 & 1 & Miss (had blk 9, replace w/ \{34,35\}) \\
11 & 1 & 0 & 0 & Miss (had blk 16, replace w/ \{0,1\}) \\
12 & 4 & 2 & 2 & Hit (blk 2: \{4,5\}) \\
\hline
\end{tabular}
\end{center}

\medskip
Hits = 6, Misses = 6

\textbf{Hit rate = 6/12 = 50\%}

\medskip
Final cache state:

\begin{center}
\begin{tabular}{|c|c|c|}
\hline
\textbf{Idx} & \textbf{Blk \#} & \textbf{Words} \\
\hline
0 & 0 & \{0, 1\} \\
1 & 17 & \{34, 35\} \\
2 & 2 & \{4, 5\} \\
3 & 3 & \{6, 7\} \\
4 & 4 & \{8, 9\} \\
5 & 5 & \{10, 11\} \\
6 & 6 & \{12, 13\} \\
7 & 7 & \{14, 15\} \\
\hline
\end{tabular}
\end{center}

\newpage
%-----------------------------------------------------------------------------
\subsection*{(c) 2-Way Set-Associative, 2-Word Blocks, 8 Blocks Total, LRU}

8 blocks / 2 ways = 4 sets. Block \# $= \lfloor \text{Addr}/2 \rfloor$, \quad Set $= \text{Block \#} \bmod 4$.

Initial state (addrs 0--15 loaded in order, so blocks 0--3 are LRU, blocks 4--7 are MRU):

\begin{center}
\begin{tabular}{|c|c|c|}
\hline
\textbf{Set} & \textbf{Way 0 (LRU)} & \textbf{Way 1 (MRU)} \\
\hline
0 & Blk 0: \{0,1\} & Blk 4: \{8,9\} \\
1 & Blk 1: \{2,3\} & Blk 5: \{10,11\} \\
2 & Blk 2: \{4,5\} & Blk 6: \{12,13\} \\
3 & Blk 3: \{6,7\} & Blk 7: \{14,15\} \\
\hline
\end{tabular}
\end{center}

\begin{center}
\small
\begin{tabular}{|c|c|c|c|c|p{5.8cm}|}
\hline
\textbf{\#} & \textbf{Addr} & \textbf{Blk} & \textbf{Set} & \textbf{H/M} & \textbf{Set After (LRU $\to$ MRU)} \\
\hline
1 & 1 & 0 & 0 & Hit & Blk4:\{8,9\}, Blk0:\{0,1\} \\
\hline
2 & 20 & 10 & 2 & Miss & Evict Blk2 $\to$ Blk6:\{12,13\}, Blk10:\{20,21\} \\
\hline
3 & 2 & 1 & 1 & Hit & Blk5:\{10,11\}, Blk1:\{2,3\} \\
\hline
4 & 3 & 1 & 1 & Hit & Blk5:\{10,11\}, Blk1:\{2,3\} \\
\hline
5 & 4 & 2 & 2 & Miss & Evict Blk6 $\to$ Blk10:\{20,21\}, Blk2:\{4,5\} \\
\hline
6 & 18 & 9 & 1 & Miss & Evict Blk5 $\to$ Blk1:\{2,3\}, Blk9:\{18,19\} \\
\hline
7 & 5 & 2 & 2 & Hit & Blk10:\{20,21\}, Blk2:\{4,5\} \\
\hline
8 & 19 & 9 & 1 & Hit & Blk1:\{2,3\}, Blk9:\{18,19\} \\
\hline
9 & 33 & 16 & 0 & Miss & Evict Blk4 $\to$ Blk0:\{0,1\}, Blk16:\{32,33\} \\
\hline
10 & 34 & 17 & 1 & Miss & Evict Blk1 $\to$ Blk9:\{18,19\}, Blk17:\{34,35\} \\
\hline
11 & 1 & 0 & 0 & Hit & Blk16:\{32,33\}, Blk0:\{0,1\} \\
\hline
12 & 4 & 2 & 2 & Hit & Blk10:\{20,21\}, Blk2:\{4,5\} \\
\hline
\end{tabular}
\end{center}

\medskip
Hits = 7, Misses = 5

\textbf{Hit rate = 7/12 $\approx$ 58.33\%}

\medskip
Final cache state:

\begin{center}
\begin{tabular}{|c|c|c|}
\hline
\textbf{Set} & \textbf{Way 0 (LRU)} & \textbf{Way 1 (MRU)} \\
\hline
0 & Blk 16: \{32,33\} & Blk 0: \{0,1\} \\
1 & Blk 9: \{18,19\} & Blk 17: \{34,35\} \\
2 & Blk 10: \{20,21\} & Blk 2: \{4,5\} \\
3 & Blk 3: \{6,7\} & Blk 7: \{14,15\} \\
\hline
\end{tabular}
\end{center}

\newpage
%=============================================================================
% PROBLEM 3
%=============================================================================
\section*{3. AMAT and CPI (20 pts)}

Given:
\begin{itemize}[nosep]
  \item \$I cache: 90\% hit rate, 1 cycle hit time, 100 cycle miss penalty
  \item \$D cache: 85\% hit rate, 2 cycle hit time, 120 cycle miss penalty
  \item Memory instructions (lw/sw) are 30\% of the instruction mix
  \item Base CPI = 1.0 (perfect memory)
\end{itemize}

\subsection*{(a) AMAT}

$\text{AMAT} = \text{Hit Time} + \text{Miss Rate} \times \text{Miss Penalty}$

\medskip
\$I cache:
\[
\text{AMAT}_I = 1 + 0.10 \times 100 = 11 \text{ cycles}
\]

\$D cache:
\[
\text{AMAT}_D = 2 + 0.15 \times 120 = 20 \text{ cycles}
\]

\subsection*{(b) CPI with Stalls}

\$I cache stall (every instruction accesses \$I):
\[
\text{Stall}_I = 0.10 \times 100 = 10 \text{ cycles/instruction}
\]

\$D cache stall (only 30\% of instructions access \$D):
\[
\text{Stall}_D = 0.30 \times 0.15 \times 120 = 5.4 \text{ cycles/instruction}
\]

Overall:
\[
\text{CPI} = 1.0 + 10 + 5.4 = \textbf{16.4}
\]

\newpage
%=============================================================================
% PROBLEM 4
%=============================================================================
\section*{4. Cache Optimization Strategies (20 pts)}

\subsection*{(a) Graph Applications (irregular access patterns)}

Graph workloads like traversing a Twitter follower graph have random, pointer-chasing access patterns with poor spatial and temporal locality.

\begin{itemize}[leftmargin=*]
  \item \textbf{Higher associativity} -- Random accesses cause lots of conflict misses. Using higher associativity (like 8-way or fully associative) helps reduce these since different addresses won't keep evicting each other from the same set.

  \item \textbf{Smaller block size} -- Spatial locality is bad here, so large cache lines just waste bandwidth fetching neighboring data that won't be used. Smaller blocks are more efficient.

  \item \textbf{Non-blocking caches} -- Graph traversals can have many independent memory accesses happening. A non-blocking cache lets multiple misses overlap so the processor isn't just sitting idle waiting for one miss at a time.

  \item \textbf{Disable stride prefetching} -- Standard hardware prefetchers assume sequential/stride patterns, which doesn't apply here. They'd just pollute the cache with wrong data, so it's better to turn them off for this kind of workload.

  \item \textbf{Victim cache} -- A small fully-associative victim cache can catch blocks that just got evicted and might be needed again soon (like revisiting a graph node), which helps with conflict misses without much extra cost.
\end{itemize}

\subsection*{(b) Streaming Applications (regular access patterns)}

Image/video processing scans through data sequentially with very predictable patterns -- great spatial locality but low temporal locality since each piece of data is usually only used once.

\begin{itemize}[leftmargin=*]
  \item \textbf{Larger block size} -- Sequential access means neighboring data \textit{will} be used, so bigger cache lines take advantage of the spatial locality and reduce the miss rate.

  \item \textbf{Hardware prefetching} -- The regular stride pattern is exactly what hardware prefetchers are designed for. They can detect the pattern and bring data into the cache ahead of time, hiding memory latency.

  \item \textbf{Lower associativity} -- Since the access pattern is predictable and doesn't cause many conflict misses, a direct-mapped or 2-way cache is fine. This gives faster hit times and uses less power.

  \item \textbf{Write buffer} -- Streaming apps that produce output (like processed video frames) generate lots of writes. A write buffer absorbs these so the processor doesn't stall waiting for each write to finish.

  \item \textbf{Cache bypassing / stream buffer} -- If the data is truly only accessed once, it can bypass the cache entirely (or go through a stream buffer) so it doesn't pollute the cache and evict other useful data.
\end{itemize}

\end{document}
